% ============================================================
%  CHƯƠNG 2: NỀN TẢNG TOÁN HỌC
%  Người viết: [Họ tên thành viên phụ trách]
% ============================================================

\section{Nền tảng toán học}
\label{sec:math_foundations}

% -----------------------------------------------------------
\subsection{Biến ngẫu nhiên và xác suất có điều kiện}
% -----------------------------------------------------------
% Gợi ý nội dung:
%   - Định nghĩa biến ngẫu nhiên liên tục, hàm mật độ xác suất
%   - Quy tắc nhân xác suất: p(x,z) = p(x|z)p(z)
%   - Định lý Bayes: p(z|x) = p(x|z)p(z) / p(x)
%   - Xác suất biên: p(x) = ∑_z p(x,z) hoặc ∫ p(x,z)dz

% TODO: Viết nội dung ở đây


% -----------------------------------------------------------
\subsection{Phân phối Gaussian đa biến}
% -----------------------------------------------------------
% Gợi ý nội dung:
%   - Định nghĩa N(x | μ, Σ), công thức mật độ
%   - Ý nghĩa của vector trung bình μ và ma trận hiệp phương sai Σ
%   - Tính chất: đường đồng mức ellipsoid, phân phối biên vẫn là Gaussian
%   - Trường hợp đặc biệt: Σ = σ²I (isotropic), Σ = diag(σ₁²,…,σ_D²)

% TODO: Viết nội dung ở đây


% -----------------------------------------------------------
\subsection{Log-likelihood và ước lượng hợp lý cực đại (MLE)}
% -----------------------------------------------------------
% Gợi ý nội dung:
%   - Hàm likelihood L(θ; X) và log-likelihood ℓ(θ)
%   - Nguyên lý MLE: tìm θ* = argmax ℓ(θ)
%   - Vì sao dùng log? (tránh underflow, biến tích thành tổng)
%   - Ví dụ: MLE cho Gaussian đơn biến

% TODO: Viết nội dung ở đây


% -----------------------------------------------------------
\subsection{Bất đẳng thức Jensen}
% -----------------------------------------------------------
% Gợi ý nội dung:
%   - Phát biểu: f hàm lồi ⟹ f(E[X]) ≤ E[f(X)]
%   - Trường hợp hàm lõm (log): E[log X] ≤ log E[X]
%   - Vai trò trong chứng minh hội tụ của EM (ELBO ≤ log p(x|θ))

% TODO: Viết nội dung ở đây


% -----------------------------------------------------------
\subsection{Kỳ vọng của biến ẩn và hàm Q}
% -----------------------------------------------------------
% Gợi ý nội dung:
%   - Dữ liệu đầy đủ (complete data) vs. dữ liệu quan sát (observed data)
%   - Định nghĩa hàm Q(θ | θ^(t)) = E_{Z|X,θ^(t)}[log p(X,Z|θ)]
%   - Đây là nền tảng lý thuyết để hiểu EM ở chương tiếp theo

% TODO: Viết nội dung ở đây
